\begin{abstract}
% v12.2, 2026-09-03. Built sentence by sentence against the 48-abstract corpus in
% paper/reference/neighbor_abstracts.md and the rules in paper/rules/02_abstract.md.
% Decisions recorded in .workspace/plans/abstract_rewrite_plan.md.
% 197 words / 9 sentences. E2 register audit passed: 0 em dashes, 0 colons,
% 0 spelled fractions, 0 assertion verbs, 0 imperatives, 1 result statistic,
% evidence as grammatical subject in the close.
% v12.2 fixes a line-wrap that split "self-correction" at its hyphen in v12.1.
Large language models are increasingly handed work their users could not do themselves.
Uncertain whether the output is right and unable to name what is wrong, users fall back on
undirected requests that carry no information about what to change. Studies of
self-correction ask whether a model can repair a response that is wrong, and find that it
can when given a specific critique and largely cannot without one. That leaves unexamined
the case where the work is already sufficient and gets revised anyway because the user
cannot tell. In this work, we run five-turn revision conversations with six models on forty
tasks across five domains, contrasting undirected requests with a single targeted critique.
Most replies contain no revision at all, with models restating the work or declining while
presenting it as compliance. Among the revisions that do change quality, 70\% lower it, and
the pattern holds across all five domains. For every model the first draft rates highest,
and blind readers prefer it to the final version only slightly more often than chance. These
results indicate that undirected revision lowers the quality of work that was already
sufficient, and that naming the fault reverses the effect.
\end{abstract}
